%% Online Resource 1 del paper DRL-vs-GP para el IJSP.
%% Mismo ciclo de compilacion que main.tex (paper/compila.bat lo cubre).
\documentclass[pdflatex,sn-apa]{sn-jnl}

\usepackage{lmodern}
\usepackage[T1]{fontenc}
\usepackage{graphicx}
\usepackage{amsmath,amssymb}
\usepackage{booktabs}

\providecommand{\doi}[1]{\href{https://doi.org/#1}{\nolinkurl{https://doi.org/#1}}}

\raggedbottom

\begin{document}

\title[Supplementary material]{Supplementary material for: Deep
Reinforcement Learning for the Interval Job Shop Scheduling Problem:
A Comparison with Genetic Programming Hyper-Heuristics across
Inference Budgets}

\author*[1]{\fnm{Hern\'an} \sur{D\'iaz Rodr\'iguez}}\email{diazhernan@uniovi.es}

\affil*[1]{\orgdiv{Department of Computer Science},
\orgname{University of Oviedo}, \orgaddress{\city{Gij\'on},
\country{Spain}}}

\abstract{This supplementary material contains three blocks
referenced from the paper: the automatic
hyperparameter-configuration campaigns and their search spaces
(Section~\ref{sup:config}), the architecture and controlled ablation
of the self-attention variant (Section~\ref{sup:attention}), and the
per-instance results on the 70-instance benchmark
(Section~\ref{sup:perinstance}). Table and equation numbers refer to
this document; all other references are to the paper.}

\maketitle

\section{The hyperparameter-configuration campaigns}
\label{sup:config}

The values of Table~5 of the paper were fixed by hand during
development, most of them at the defaults customary in the PPO
literature. To assess whether automatic tuning would improve on them,
two configuration studies were run with
irace~\citep{LopezIbanez2016}. The first raced $330$ experiments
($309$ used) over the 8-parameter space of Table~\ref{sup:tab-irace}
at reduced fidelity: one training instance, $300$ episodes,
${\sim}8$ minutes per evaluation. Its elites agreed on a more
aggressive update regime and kept minibatch size and update period at
their defaults, but at the full $1000$-episode operating budget the
best of them scored $14.5\%$ mean RE against $13.4\%$ for the
defaults: the low-fidelity optimum did not transfer.

\begin{table}[t]
\centering\small
\caption{Search space of the two irace campaigns. Real parameters were
sampled on a logarithmic scale, the rest from the listed values.}
\label{sup:tab-irace}
\setlength{\tabcolsep}{5pt}
\begin{tabular}{llll}
\toprule
Parameter & Default & Campaign 1 (330 exp.) & Campaign 2 (300 exp.) \\
\midrule
learning rate & $3\cdot10^{-4}$ & $[10^{-4},\,10^{-3}]$ &
  $[10^{-4},\,10^{-3}]$ \\
entropy coefficient & $0.01$ & $[0.003,\,0.03]$ & $[0.003,\,0.03]$ \\
clipping range & $0.2$ & $0.1,\ 0.2,\ 0.3$ & $0.1,\ 0.2,\ 0.3$ \\
epochs per update & $4$ & $2,\ 4,\ 8$ & $4,\ 8$ \\
GAE $\lambda$ & $0.95$ & $0.90,\ 0.95,\ 0.98$ & $0.90,\ 0.95,\ 0.98$ \\
hidden width & $128$ & $64,\ 128,\ 256$ & $128,\ 256$ \\
minibatch size & $256$ & $128,\ 256,\ 512$ & fixed \\
update period & $4$ ep. & $2,\ 4,\ 8$ ep. & fixed \\
\bottomrule
\end{tabular}
\end{table}

The second campaign raced $295$ experiments at operating fidelity,
with $1000$ episodes on two instances per evaluation, over the
reduced $6$-parameter space of Table~\ref{sup:tab-irace}, seeded with
the defaults and the first campaign's elites. It converged to a
region distinct from the defaults and eliminated the seeded default
during racing; yet its winner, retrained on the four training
instances over three seeds, again failed to improve: $14.73\%$
against $13.52\%$ over the six validation instances.
A third campaign raced the reward instead of the optimizer: the six
weights of Eq.~(5) of the paper, each free in $[0,1]$, over $299$
experiments at operating fidelity. Its winner did not improve the
deployed weights either, $14.04\%$ against $13.55\%$ under the same
confirmation protocol (Wilcoxon signed-rank, $p=0.21$). Two
by-products of that campaign are more informative than its verdict. A
seeded terminal-only vector $(1,0,0,0,0,0)$ was eliminated in the
first round, so the dense shaping is not dispensable; and four elites
survived the full race differing by as much as $0.48$ in a single
weight, with a Kendall agreement across seeds of $W=0.02$--$0.05$ on
which is better, so at this budget the objective is flat in the
reward weights relative to the noise between training runs.

Automatic configuration therefore does not improve on the hand-set
values for this problem, even at the operating budget, nor do raced
reward weights improve on the hand-set ones, within the resolution of
the protocol. That resolution can be quantified: with a standard
deviation of $0.47$ points across seed means at $1000$ episodes and
three seeds per arm, the minimum difference a paired confirmation
detects at conventional power is approximately one point of RE. The
three campaigns therefore exclude improvements larger than that
magnitude, not the existence of smaller ones. The same arithmetic
explains their common signature, the racing optimum never matching
the confirmation optimum: single-run outcomes are too noisy to
separate near-equivalent configurations, and raising the resolution
by repeating each evaluation reduces the number of configurations a
fixed budget can race.

\section{The self-attention variant}
\label{sup:attention}

Pooling summarizes the candidate set by two statistics, against which
every candidate is then scored. This construction cannot express a
comparison between two \emph{particular} candidates: whether
operation $i$ should take the machine now may depend on which other
eligible operation would occupy it next, and the mean embedding does
not identify that candidate. The embeddings of all
$|\mathcal{E}|$ candidates are available at the decision step, so
what the architecture lacks is not the information but a path along
which one candidate can condition another.
\emph{Self-attention}~\citep{Vaswani2017} supplies such a path and
is the component most commonly adopted for this purpose in the
neural combinatorial optimization literature. In an attention layer
every candidate emits a \emph{query}, a \emph{key} and a
\emph{value}; the similarity between one candidate's query and
another's key decides how much of that other's value is added into
the first's embedding. Each candidate is thereby rewritten as a
content-weighted reading of all the others, with the weights computed
from the embeddings themselves rather than fixed in advance.

A stack of $B$ pre-layer-normalization (pre-LN) transformer blocks
transforms the embeddings between the embedding and pooling stages of
the base architecture (Figure~1 of the paper). Writing
$\Phi\in\mathbb{R}^{|\mathcal{E}|\times h}$ for the stacked
embeddings, each block computes
\begin{equation}
\label{sup:eq-attnblock}
Z=\Phi+\mathrm{MHA}\big(\mathrm{LN}(\Phi)\big),\qquad
\Phi'=Z+\mathrm{FFN}\big(\mathrm{LN}(Z)\big),
\end{equation}
where $\mathrm{MHA}$ is multi-head scaled dot-product attention ($4$
heads of dimension $32$, no dropout), $\mathrm{FFN}$ is a two-layer
perceptron ($128\to256\to128$, ReLU) applied to each row, and
$\mathrm{LN}$ is LayerNorm. Both additions are \emph{residual}, that
is, each block adds to its input rather than replacing it, so an
untrained block is close to the identity and the stack starts near
the base model; the pre-LN arrangement is preferred over the original
post-LN one for this property~\citep{Xiong2020}. No positional
encoding is used, which keeps self-attention
permutation-equivariant, so the pooled context remains
permutation-invariant and the size invariance of the base
architecture is preserved. Padding slots are masked out of the
attention weights, and unit tests verify that masked-padded forward
passes are numerically identical to unpadded ones.

The variant is a configuration option on the same code path, and
$B{=}0$ reproduces the base network exactly, so the comparison below
differs in this component and in nothing else. A single block carries
$132{,}480$ parameters in the query, key and value projections, the
output projection, two layer normalizations and the feed-forward
network, so the $B{=}2$ variant evaluated here grows the network from
$1.2\times10^{5}$ to $3.9\times10^{5}$ parameters, a factor of
$3.2$, and multiplies the wall-clock cost of an episode by $2.2$.

\paragraph{The controlled ablation.} The variant inserts two residual
attention blocks between the shared per-operation encoder and the
masked pooling, which it leaves in place; everything else is held
fixed, including the training and evaluation protocols, the same four
training instances, the same six validation instances, the same seeds
and the hyperparameters of Table~5 of the paper. The two arms are
therefore matched instance by instance and seed by seed, and analysed
with the instance as the experimental unit, at both training budgets
employed in the study: over three seeds at $300$ episodes and over
ten at $1000$, the operating budget.

The variant yields no improvement at either budget and a significant
degradation at the operating one. At $300$ episodes it is $1.03$
points of RE inferior (95\% CI $0.01$ to $2.06$), degraded on 5 of
the 6 instances, a difference that does not reach significance
(Wilcoxon signed-rank, $p=0.063$); at $1000$ episodes the deficit is
$1.06$ points ($0.46$ to $1.65$), degraded on all six instances and
in every one of the ten seeds ($p=0.031$, the smallest value six
instances admit). Table~\ref{sup:tab-attn} reports the corresponding
means and best seeds.

\begin{table}[!htb]
\centering\small
\caption{Attention ablation: mean [best seed] RE (\%) on the
validation set, over three seeds at $300$ episodes and ten at
$1000$.}
\label{sup:tab-attn}
\begin{tabular}{lcc}
\toprule
Configuration & 300 eps & 1000 eps \\
\midrule
Deep Sets & \textbf{16.6} [16.2] & \textbf{13.8} [13.1] \\
+ attention ($B{=}2$) & 17.6 [17.2] & 14.9 [14.2] \\
\bottomrule
\end{tabular}
\end{table}

The representation accounts for part of this outcome. One of the
per-operation features is itself relational with respect to the
eligible set: slack measures the candidate's earliest start against
the earliest among all candidates, so the comparison an attention
layer would have to learn along that axis is already available as a
precomputed input. A second, machine congestion, is relational with
respect to the machines rather than the candidates, contrasting a
machine's remaining load with the average over machines. What remains
are the pairwise relations these aggregates do not summarize, and at
the budgets tested they do not amortize their cost; larger training
budgets or more heterogeneous instance distributions may be required
for such structure to become profitable.

\section{Per-instance results}
\label{sup:perinstance}

Table~\ref{sup:tab-perinstance} lists, for each of the 70 interval
Taillard instances, the crisp lower bound and the RE (\%) of the
earliest-start rule, of Giffler--Thompson with MWKR tie-breaking, of
the evolved rule of~\citet{DiazGP2026}, and of the policy at one pass
and at 1024 samples, the artifact each study selects as in Table~8 of
the paper. The ten $20{\times}15$ instances TA11--TA20 are the
training and validation set.

\begin{table}[h]
\centering\scriptsize
\caption{Per-instance RE (\%) of the constructive methods and of the
policy. LB is the crisp Taillard lower bound; GP the evolved rule;
Pol.\ the policy at one pass and Pol.$^{\ast}$ at 1024 samples.}
\label{sup:tab-perinstance}
\renewcommand{\arraystretch}{0.92}
\setlength{\tabcolsep}{3pt}
\begin{tabular}{lrrrrrr@{\hskip 0.7em}|@{\hskip 0.7em}lrrrrrr}
\toprule
Inst. & LB & EST & G\&T & GP & Pol. & Pol.$^{\ast}$ &
Inst. & LB & EST & G\&T & GP & Pol. & Pol.$^{\ast}$ \\
\midrule
TA1 & 1231 & 49.9 & 26.1 & 19.3 & 17.2 & 12.3 & TA36 & 1819 & 41.1 & 42.8 & 23.8 & 19.2 & 13.4 \\
TA2 & 1244 & 24.9 & 27.8 & 14.2 & 14.8 & 7.9 & TA37 & 1771 & 43.4 & 35.6 & 18.4 & 22.4 & 17.7 \\
TA3 & 1218 & 33.0 & 33.2 & 14.1 & 14.0 & 8.2 & TA38 & 1673 & 45.0 & 28.0 & 20.1 & 20.1 & 15.7 \\
TA4 & 1175 & 32.0 & 27.7 & 13.0 & 18.5 & 11.3 & TA39 & 1795 & 46.6 & 28.9 & 22.4 & 22.3 & 14.7 \\
TA5 & 1224 & 27.2 & 33.0 & 14.7 & 10.6 & 7.8 & TA40 & 1651 & 46.9 & 42.3 & 26.1 & 27.8 & 19.4 \\
TA6 & 1238 & 26.8 & 18.2 & 13.7 & 14.8 & 7.7 & TA41 & 1906 & 44.6 & 45.0 & 33.4 & 27.8 & 24.3 \\
TA7 & 1227 & 30.7 & 26.1 & 17.7 & 15.2 & 9.1 & TA42 & 1884 & 71.1 & 38.1 & 24.3 & 23.7 & 20.3 \\
TA8 & 1217 & 23.7 & 25.0 & 14.8 & 17.3 & 9.8 & TA43 & 1809 & 68.7 & 39.8 & 22.3 & 27.4 & 20.7 \\
TA9 & 1274 & 37.7 & 28.2 & 20.0 & 12.9 & 11.1 & TA44 & 1948 & 52.1 & 34.7 & 25.8 & 27.4 & 17.9 \\
TA10 & 1241 & 37.0 & 22.1 & 15.4 & 15.1 & 7.6 & TA45 & 1997 & 50.2 & 33.7 & 14.9 & 16.7 & 13.0 \\
TA11 & 1357 & 56.2 & 30.7 & 17.5 & 19.8 & 12.4 & TA46 & 1957 & 60.8 & 31.6 & 22.4 & 21.7 & 19.6 \\
TA12 & 1367 & 48.8 & 34.6 & 17.6 & 13.0 & 11.1 & TA47 & 1807 & 46.8 & 30.1 & 25.6 & 24.4 & 19.8 \\
TA13 & 1342 & 60.2 & 26.9 & 13.3 & 19.7 & 12.0 & TA48 & 1912 & 59.1 & 33.7 & 21.1 & 22.8 & 19.3 \\
TA14 & 1345 & 39.5 & 29.7 & 13.2 & 17.7 & 9.0 & TA49 & 1931 & 51.0 & 35.3 & 20.1 & 23.3 & 18.7 \\
TA15 & 1339 & 48.1 & 30.6 & 16.7 & 17.0 & 11.6 & TA50 & 1833 & 61.4 & 46.2 & 31.3 & 32.2 & 25.6 \\
TA16 & 1360 & 45.2 & 34.0 & 15.6 & 15.9 & 9.3 & TA51 & 2760 & 33.8 & 34.7 & 23.2 & 24.3 & 13.4 \\
TA17 & 1462 & 58.9 & 17.7 & 13.9 & 14.5 & 9.9 & TA52 & 2756 & 38.0 & 26.5 & 12.7 & 19.0 & 12.7 \\
TA18 & 1377 & 43.8 & 31.8 & 23.6 & 20.9 & 14.1 & TA53 & 2717 & 34.1 & 18.3 & 11.8 & 13.5 & 9.7 \\
TA19 & 1332 & 45.8 & 29.3 & 20.0 & 14.8 & 10.0 & TA54 & 2839 & 20.9 & 16.1 & 8.2 & 9.9 & 5.2 \\
TA20 & 1348 & 45.7 & 24.1 & 15.1 & 20.7 & 12.8 & TA55 & 2679 & 31.7 & 29.2 & 15.5 & 20.0 & 14.4 \\
TA21 & 1642 & 33.4 & 28.0 & 20.2 & 17.6 & 11.2 & TA56 & 2781 & 34.8 & 19.8 & 10.9 & 10.9 & 9.2 \\
TA22 & 1561 & 44.2 & 30.8 & 17.4 & 17.7 & 15.4 & TA57 & 2943 & 32.9 & 20.0 & 9.6 & 15.3 & 10.1 \\
TA23 & 1518 & 42.8 & 35.3 & 17.8 & 20.2 & 12.2 & TA58 & 2885 & 37.4 & 26.3 & 13.6 & 15.3 & 10.1 \\
TA24 & 1644 & 56.3 & 23.8 & 14.1 & 13.4 & 11.0 & TA59 & 2655 & 26.6 & 28.2 & 16.7 & 21.5 & 15.0 \\
TA25 & 1558 & 45.8 & 33.8 & 16.6 & 18.3 & 11.0 & TA60 & 2723 & 41.7 & 19.8 & 15.5 & 11.1 & 9.0 \\
TA26 & 1591 & 42.8 & 31.2 & 26.9 & 18.8 & 13.1 & TA61 & 2868 & 45.7 & 25.4 & 15.4 & 18.2 & 15.2 \\
TA27 & 1652 & 50.0 & 31.1 & 21.4 & 18.3 & 13.2 & TA62 & 2869 & 43.0 & 27.6 & 17.2 & 20.2 & 15.9 \\
TA28 & 1603 & 36.8 & 26.8 & 12.4 & 15.6 & 9.3 & TA63 & 2755 & 38.8 & 23.8 & 13.4 & 17.8 & 12.3 \\
TA29 & 1583 & 28.8 & 32.0 & 17.2 & 14.9 & 9.8 & TA64 & 2702 & 38.2 & 28.4 & 12.2 & 13.2 & 12.2 \\
TA30 & 1528 & 34.2 & 36.1 & 17.0 & 21.9 & 15.6 & TA65 & 2725 & 42.9 & 30.0 & 20.5 & 20.8 & 17.3 \\
TA31 & 1764 & 30.7 & 34.4 & 24.5 & 20.5 & 15.9 & TA66 & 2845 & 41.8 & 28.0 & 11.9 & 15.7 & 13.2 \\
TA32 & 1774 & 42.7 & 33.3 & 19.0 & 25.1 & 18.1 & TA67 & 2825 & 31.7 & 26.1 & 15.2 & 15.9 & 13.0 \\
TA33 & 1788 & 51.3 & 34.8 & 27.5 & 30.7 & 21.0 & TA68 & 2784 & 50.1 & 19.1 & 10.6 & 15.4 & 10.6 \\
TA34 & 1828 & 45.9 & 31.8 & 20.7 & 18.2 & 16.0 & TA69 & 3071 & 35.3 & 22.8 & 12.9 & 15.4 & 11.6 \\
TA35 & 2007 & 30.3 & 24.1 & 8.4 & 7.9 & 4.1 & TA70 & 2995 & 42.9 & 23.4 & 16.7 & 18.4 & 15.3 \\
\bottomrule
\end{tabular}
\end{table}

% el estilo lo fija la opcion de clase (sn-apa -> sn-apacite.bst)
\bibliography{refs}

\end{document}
